
 
% ==========================================
% CHƯƠNG 7: KẾT LUẬN VÀ HƯỚNG PHÁT TRIỂN
% ==========================================
\chapter{KẾT LUẬN VÀ HƯỚNG PHÁT TRIỂN}
 
\section{Kết luận}
 
Đề tài đã nghiên cứu và xây dựng thành công một hệ thống dịch hai chiều toàn diện giữa Ngôn ngữ Ký hiệu Việt Nam (VSL) và tiếng Việt, góp phần thiết thực vào mục tiêu xóa bỏ rào cản giao tiếp cho cộng đồng người khiếm thính. Các kết quả trọng tâm đã đạt được bao gồm:
 
\begin{itemize}
    \item Làm chủ công nghệ nhận diện đa thành phần với feature vector 178 chiều (Pose + Hands + Finger Curl + Emotion), đặc biệt là việc tích hợp thành công 7 chiều Emotion Encoding -- yếu tố quyết định giúp phân biệt các ký hiệu đồng hình trong VSL.
 
    \item Xây dựng hoàn chỉnh hai luồng xử lý: Chiều VSL $\rightarrow$ VN với mô hình Transformer + Attention + PhoBERT, và chiều VN $\rightarrow$ VSL với NLP Pipeline sử dụng Underthesea, Sliding Window và cơ sở dữ liệu video ký hiệu.
 
    \item Đạt hiệu suất thời gian thực nhờ bộ đệm cửa sổ trượt, xử lý bất đồng bộ MediaPipe và Stability Filter với confidence scoring tổng hợp.
 
    \item Đóng góp quy trình thu thập và tăng cường dữ liệu VSL có cấu trúc, kết hợp 7 trạng thái cảm xúc, tạo ra bộ dữ liệu đa dạng và giàu ngữ nghĩa hơn các công trình trước đây.
\end{itemize}
 
\section{Hướng phát triển}
 
Nhóm nghiên cứu định hướng mở rộng và hoàn thiện hệ thống theo các trục chính:
 
\begin{itemize}
    \item \textbf{Mở rộng quy mô bộ dữ liệu:} Hướng tới bộ dữ liệu VSL trên 10.000 từ vựng, bao gồm cụm từ chuyên ngành, phương ngữ vùng miền và đa dạng hơn về người thực hiện, điều kiện thu thập.
 
    \item \textbf{Nâng cấp kiến trúc mô hình:} Nghiên cứu áp dụng kiến trúc Transformer thuần túy (Vision Transformer, Temporal Transformer) hoặc Multimodal Learning để khai thác tốt hơn mối quan hệ giữa cử chỉ, cảm xúc và ngữ cảnh.
 
    \item \textbf{Tối ưu hóa cho thiết bị nhúng:} Áp dụng các kỹ thuật model compression (quantization, pruning) để triển khai mô hình trực tiếp trên kính thông minh hoặc điện thoại di động, giảm phụ thuộc vào kết nối server.
 
    \item \textbf{Tích hợp công nghệ mới:} Phát triển tính năng hội thoại video thời gian thực với phụ đề tự động, tích hợp Thực tế ảo tăng cường (AR) để học ký hiệu tương tác, và Text-to-Speech đa ngôn ngữ để hỗ trợ giao tiếp quốc tế.
 
    \item \textbf{Triển khai xã hội:} Hợp tác với các trường chuyên biệt, bệnh viện và cơ quan hành chính để đưa hệ thống vào ứng dụng thực tiễn, đóng góp trực tiếp cho mục tiêu bình đẳng giao tiếp và hòa nhập xã hội cho người khiếm thính Việt Nam.
\end{itemize}
 
Hệ thống được kỳ vọng sẽ không chỉ là một công cụ dịch thuật đơn thuần mà còn trở thành nền tảng công nghệ thiết yếu, góp phần kiến tạo một xã hội hòa nhập, bình đẳng và không rào cản giao tiếp cho hơn 2,5 triệu người khiếm thính tại Việt Nam.